\centerline{\bf \Huge 10 класс}

\centerline{\bf \Large Первый день}

\problem{10.1}
Найдите НОД всех чисел вида $n^{13}-n$, где $n$ - натуральное число.

\problem{10.2}
Дана равнобокая трапеция $A B C D$ с основаниями $B C$ и $A D$. Окружность $\omega$ проходит через вершины $B$ и $C$ и вторично пересекает сторону $A B$ и диагональ $B D$ в точках $X$ и $Y$ соответственно. Касательная, проведенная к окружности $\omega$ в точке $C$, пересекает луч $A D$ в точке $Z$. Докажите, что точки $X, Y$ и $Z$ лежат на одной прямой.

\problem{10.3}
Клетки таблицы $2 \times 2025$ надо заполнить числами (в каждую клетку вписать ровно одно число) по следующим правилам. В верхней строке должны стоять 2025 действительных чисел, среди которых нет двух равных, а в нижней строке должны стоять те же 2025 чисел, но в другом порядке. В каждом из 2025 столбцов должны быть записаны два разных числа, причём сумма этих двух чисел должна быть рациональным числом. Какое наибольшее количество иррациональных чисел могло быть в первой строке таблицы?

\problem{10.4}
Бесконечная последовательность ненулевых чисел $a_1, a_2, a_3, \ldots$ такова, что при всех натуральных $n \geqslant 2025$ число $a_{n+1}$ является наименьшим корнем многочлена

$$
P_n(x)=x^{2 n}+a_1 x^{2 n-2}+a_2 x^{2 n-4}+\ldots+a_n .
$$


Докажите, что существует такое $N$, что в бесконечной последовательности $a_N, a_{N+1}, \\a_{N+2}, \ldots$ каждый член меньше предыдущего.

\problem{10.5}
Из каждого угла квадрата $300 \times 300$ вырезали по квадрату $100 \times 100$, а оставшиеся 50000 клеток раскрасили в черный и белый цвета так, чтобы никакой квадрат $2 \times$ 2 не был бы раскрашен в шахматном порядке. Найдите максимально возможное количество соседних пар клеток, окрашенных в разные цвета.